\subsection{Drei geometrische Lösungstypen}
\label{subsec:linear_systems_solution_sets_and_geometry}

Die Lösungsmenge eines LGS kann sehr unterschiedlich aussehen. Wir
untersuchen die möglichen Fälle zunächst an $2\times2$-LGS.

In zwei Dimensionen beschreibt jede lineare Gleichung eine Gerade in der
$x_1$-$x_2$-Ebene. Ein Lösungstupel $(x_1,x_2)$ entspricht einem gemeinsamen
Punkt der Geraden.
\begin{itemize}[leftmargin=1.6em,itemsep=1.2em]
  \item
    \begin{minipage}[t]{\linewidth}
    \textbf{LGS hat eine eindeutige Lösung}\par\smallskip
    \noindent
    \begin{minipage}[c]{0.64\linewidth}
      Wir betrachten das LGS
      \[
        \begin{aligned}
          \textcolor{THMBlue}{x_1+x_2}&\textcolor{THMBlue}{=2},\\
          \textcolor{black!60}{x_1-x_2}&\textcolor{black!60}{=0}.
        \end{aligned}
        \quad\Longleftrightarrow\quad
        \left[
          \begin{array}{rr|r}
            \textcolor{THMBlue}{1}&\textcolor{THMBlue}{1}&\textcolor{THMBlue}{2}\\
            \textcolor{black!60}{1}&\textcolor{black!60}{-1}&\textcolor{black!60}{0}
          \end{array}
        \right]
      \]
      Durch Addition beider Gleichungen folgt $2x_1=2$ und damit $x_1=1$.
      Einsetzen in eine der Gleichungen ergibt $x_2=1$.

      \smallskip
      Geometrische Betrachtung: Die Geraden besitzen genau einen
      Schnittpunkt $S=(1,1)$. Somit ist
      \[
        \mathcal{L}=\Cloze{5cm}{\{(1,1)\}}.
      \]
    \end{minipage}\hfill
    \begin{minipage}[c]{0.32\linewidth}
      \centering
      \CachedTikzFigure{content/resources/chapters/linear_systems/tikz/unique_solution_lines.tex}
    \end{minipage}
    \end{minipage}

  \item
    \begin{minipage}[t]{\linewidth}
    \textbf{LGS hat keine Lösung}\par\smallskip
    \noindent
    \begin{minipage}[c]{0.64\linewidth}
      Wir betrachten das LGS
      \[
        \begin{aligned}
          \textcolor{THMBlue}{x_1+x_2}&\textcolor{THMBlue}{=2},\\
          \textcolor{black!60}{x_1+x_2}&\textcolor{black!60}{=1}.
        \end{aligned}
        \quad\Longleftrightarrow\quad
        \left[
          \begin{array}{rr|r}
            \textcolor{THMBlue}{1}&\textcolor{THMBlue}{1}&\textcolor{THMBlue}{2}\\
            \textcolor{black!60}{1}&\textcolor{black!60}{1}&\textcolor{black!60}{1}
          \end{array}
        \right]
      \]
      Die Gleichungen widersprechen sich gegenseitig und können daher
      nicht gleichzeitig erfüllt sein.

      \smallskip
      Geometrische Betrachtung: Die Geraden sind parallel und besitzen
      keinen Schnittpunkt. Daher gilt
      \[
        \mathcal{L}=\Cloze{5cm}{\emptyset}.
      \]
    \end{minipage}\hfill
    \begin{minipage}[c]{0.32\linewidth}
      \centering
      \CachedTikzFigure{content/resources/chapters/linear_systems/tikz/no_solution_lines.tex}
    \end{minipage}
    \end{minipage}

  \item
    \begin{minipage}[t]{\linewidth}
    \textbf{LGS hat unendlich viele Lösungen}\par\smallskip
    \noindent
    \begin{minipage}[c]{0.64\linewidth}
      Wir betrachten das LGS
      \[
        \begin{aligned}
          \textcolor{THMBlue}{x_1+x_2}&\textcolor{THMBlue}{=2},\\
          \textcolor{olive!80!black}{2x_1+2x_2}&\textcolor{olive!80!black}{=4}.
        \end{aligned}
        \quad\Longleftrightarrow\quad
        \left[
          \begin{array}{rr|r}
            \textcolor{THMBlue}{1}&\textcolor{THMBlue}{1}&\textcolor{THMBlue}{2}\\
            \textcolor{olive!80!black}{2}&\textcolor{olive!80!black}{2}&\textcolor{olive!80!black}{4}
          \end{array}
        \right]
      \]
      Die zweite Gleichung ist das Doppelte der ersten und stellt daher
      keine zusätzliche Bedingung dar. Jedes Tupel,
      das die erste Gleichung erfüllt, erfüllt auch die zweite.

      \smallskip
      Geometrische Betrachtung: Die Geraden sind identisch und besitzen
      unendlich viele Schnittpunkte. Das bedeutet nicht, dass jedes Tupel aus
      $\R^2$ eine Lösung ist: Lösungen sind genau die Punkte auf der gemeinsamen
      Geraden. Daher gilt
      \[
        \mathcal{L}=\Cloze{7.5cm}{\{(x_1,x_2)\in\R^2\mid x_1+x_2=2\}}.
      \]
    \end{minipage}\hfill
    \begin{minipage}[c]{0.32\linewidth}
      \centering
      \CachedTikzFigure{content/resources/chapters/linear_systems/tikz/infinitely_many_solutions_lines.tex}
    \end{minipage}
    \end{minipage}
\end{itemize}

\vspace{\baselineskip}

Die geometrische Betrachtung macht anschaulich, dass lediglich diese drei Fälle
auftreten können: Zwei Geraden schneiden sich, verlaufen parallel oder liegen
aufeinander.

Die in $2\mathrm{D}$ gewonnenen Erkenntnisse lassen sich auf beliebige
$m\times n$-LGS verallgemeinern.
\begin{TheoremBox}{Anzahl der Lösungen eines LGS}{linear_systems_three_solution_types}
  Jedes $m\times n$-LGS über $\R$ besitzt genau einen der drei Lösungstypen:
  \SolutionTypeCases
  Insbesondere kann ein LGS weder genau zwei noch allgemein endlich viele
  Lösungen größer als eins besitzen.
\end{TheoremBox}
